% Cote 119: offprints, reports, courses, CV and letters (typed and printed part): archive cote 119, site page 268 (PDF page 269).
% typescript. Draft compared: none.
% AI-transcribed by Claude Opus 5.5 (claude-opus-5-5, adaptive thinking, effort high), Message Batches API, batch r119_full (msgbatch_01DwZNm5Y7YsKBWyAN4sFs8w), 2026-11-13; not yet checked by eye.
\scanpage{119/268}
\mspage{15}

``décrocher'' une thèse ancien style), stimulés par une option nouvelle qui
s'offre à eux, porteraient à leur travail d'enseignement un intérêt plus actif
et une ardeur nouvelle, dont les étudiants seraient les premiers bénéficiaires,
ardeur qui de plus permettrait à ces enseignants d'atteindre à la qualification
voulue pour devenir maîtres de conférences. Il serait sans doute malaisé,
même à de plus compétents que moi en la matière, de faire aucun pronostic sur
la proportion parmi les assistants et maîtres-assistants qui seraient en mesure
de saisir l'opportunité de promotion professionnelle qui leur serait ainsi
offerte. Il semble probable en tout cas que nombreux également seront ceux
qui n'auront pas les moyens intellectuels nécessaires, et que pendant de longues
années et jusqu'à leur retraite, ces derniers constitueront, parmi les jeunes
promotions d'assistants et de maîtres-assistants faisant leur apprentissage
du métier de professeur, une sorte d'arrière-garde de maîtres-assistants à
vie, qui auront leur rôle à jouer dans l'inévitable transition entre l'Univer-%
sité d'hier et celle de demain.

7. Dans l'hypothèse où des conceptions dans l'esprit de celles esquissées ici
finissent par prévaloir au cours des prochaines années (ce qui sans doute ne
pourra se faire que si la contestation étudiante conserve sa vigueur, sans se
reposer sur les lauriers des avantages acquis), il faudrait revoir alors des
régimes de transition sans doute assez longs, caractérisés par la prédominance
numérique parmi les enseignants d'un bon nombre d'assistants\hand{,}et de maîtres- \margin{x}
assistants plus ou moins ``immobiles''\hand{,}et parmi les professeurs et maîtres de
conférences, de ceux qui se seront définis ``chercheurs'' plutôt \hand{\unsure{$\wedge$}}qu\over{'}{e} \hand{``}enseignants\hand{''}.
Parmi les nouveaux postes créés ou devenus vacants, il faudrait donc prévoir
qu'une majorité devraient être attribués à des maîtres de conférences mention
``enseignement'', pour arriver dans un temps X à l'état d'équilibre relatif du
``régime de croisière'' d'un tiers \over{\illegible{1}}{du} moins \over{\illegible{1}}{de} maîtres de conférences mention \margin{x x}
``enseignement''. Une telle politique universitaire semblerait conforme, sur
le plan global, à la proportion relativement réduite des chercheurs de haut
niveau parmi les jeunes mathématiciens qu'intéresse la carrière universitaire.
Ellé demandera cependant, de la part des départements de Mathématique à très
fort pouvoir d'attraction, tels ceux de Paris et d'Orsay, une volonté déli-%
bérée de résister à la tentation de profiter de tous les postes disponibles

\margin{268}
